\SourceLeaf{395}%
\[
\begin{array}{l}
M\\
M-2(a-r\cos A)\sin(A''-A')\\
M-2(a'-r\cos A')\sin(A-A'')\\
M-2(a''-r\cos A'')\sin(A'-A)
\end{array}
\]
Man hat also offenbar eine Gleichung von der Form
\[
\gamma+\delta r+\varepsilon rr+\zeta r^3=\theta\theta-uu=4\alpha\alpha\beta\beta
\]
wo \(\gamma,\delta,\varepsilon,\zeta\) gegebene Grössen sind, und dann wird \(r\) durch die Bedingung des Maximums offenbar aus folgender quadratischen Gleichung zu bestimmen sein
\[
\delta+2\varepsilon r+3\zeta rr=0
\]
\par
Noch leichter findet man die Coëfficienten dieser Gleichung durch folgende Betrachtung. Da das vierfache Product aus den Quadraten der halben grossen und der halben kleinen Axe einer jeden Ellipse, welche die vier Seiten des Vierecks berührt und deren Mittelpunkt zur Abscisse \(r\) hat, allgemein
\[
=\gamma+\delta r+\varepsilon rr+\zeta r^3
\]
wird, so muss dieser Ausdruck nothwendig \(=0\) werden, wenn man für \(r\) einen Werth substituirt, welcher einer der drei oben betrachteten verschwindenden Ellipsen entspricht. Diese drei Werthe sind die Abstände der beiden Halbirungspunkte der Diagonalen des Vierecks und des Halbirungspunktes der geraden Linie, welche die Durchschnitte der beiden Seiten-Paare des Vierecks verbinden, von dem Anfangspunkte der Abscissen. Ich bezeichne diese drei Punkte durch \(C,D,E\), und ihre Abscissen durch \(c,d,e\), so muss offenbar
\[
r^3+\tfrac{\varepsilon}{\zeta}rr+\tfrac{\delta}{\zeta}r+\tfrac{\gamma}{\zeta}
\]
mit dem Producte \((r-c)(r-d)(r-e)\) identisch sein; folglich ist die obige quadratische Gleichung
\[
3rr-2r(c+d+e)+cd+ce+de=0
\]
deren Wurzeln
\[
\tfrac{c+d+e}{3}+\tfrac13\BareRoot{(cc+dd+ee+cd+ce+de)}
\]
